\documentclass[11pt]{article}

\usepackage[margin=1in]{geometry}
\usepackage{booktabs}
\usepackage{longtable}
\usepackage{array}
\usepackage{xcolor}
\usepackage{enumitem}
\usepackage[hidelinks]{hyperref}
\usepackage{titlesec}

% \col{...} typesets a database column name verbatim (handles underscores).
\newcommand{\col}[1]{\texttt{\detokenize{#1}}}
\newcommand{\status}[1]{\textsf{#1}}

\titleformat{\section}{\large\bfseries}{\thesection}{0.7em}{}
\setlist{noitemsep,topsep=3pt}

\title{\bfseries A Machine-Mined Reference Database of\\
CASSCF-Family Calculations on Transition-Metal and\\
Organic Complexes: Methodology}
\author{\texttt{claude-casscf-data\_mining}}
\date{Compiled \today}

\begin{document}
\maketitle

\begin{abstract}
\noindent
This document describes the construction of \col{literature.csv}, a
structured reference database of multireference quantum-chemical
calculations---CASSCF, RASSCF, CASPT2, RASPT2, NEVPT2, DMRG, and MRCI---%
extracted from the published open-access literature. The database is built
by a semi-automated pipeline in which a Python package performs the
mechanical work (bibliographic search, PDF retrieval, and text extraction)
and a large-language-model agent performs the reading and judgment
(deciding whether a paper reports a usable calculation and transcribing the
relevant quantities into a fixed schema). The design goal is a
\emph{faithful, append-only, auditable} corpus in which every stored value
is traceable to a statement in a specific paper, so that it can serve as a
reference set for an active-space decision agent. As of the current
snapshot the database contains \textbf{667} calculation records drawn from
\textbf{403} distinct publications spanning \textbf{1985--2026}.
\end{abstract}

\section{Motivation and scope}

Choosing a complete active space (CAS) for a multireference calculation is
an expert task: the number of active electrons and orbitals, and the
chemical identity of those orbitals, must be tailored to the electronic
structure of the target molecule. A database of active spaces that
practitioners have actually published---keyed to the compound, the metal
center, the method, and the resulting state energetics---is therefore a
useful reference for guiding new calculations. This database is the
literature-mining companion to a separate active-space \emph{decision
agent}; the two share an identical schema so that mined records can be
merged directly into the decision agent's reference set.

A record is \emph{in scope} when a paper reports an actual multireference
calculation (CASSCF / RASSCF / CASPT2 / RASPT2 / NEVPT2 / MRCI / DMRG /
selected-CI-SCF) on a \emph{specific chemical compound}, with at least the
active-space size (number of electrons and orbitals) stated. Both
transition-metal complexes and metal-free organic molecules are included.
Explicitly \emph{out of scope}: review articles that merely cite others'
numbers, algorithm/method papers whose only test cases are toy systems
(diatomics or model Hamiltonians used to benchmark a code), and any paper
in which the active space is never specified.

\section{Pipeline overview}

The corpus is produced one paper at a time by a six-stage loop. Each paper
carries a \col{status} that advances monotonically through the stages and is
recorded in \col{papers_index.csv} (one row per candidate paper). A strict
division of labor is maintained: the Python package does only mechanical,
reproducible work; every act of chemical judgment is performed by the agent
and logged in prose.

\begin{enumerate}
\item \textbf{Search.} The \href{https://openalex.org}{OpenAlex} API is
  queried with an open-access filter for CASSCF-related terms. New
  candidates are appended to the index (\status{candidate}), deduplicated by
  DOI; a paper already present under any status is never re-added. Queries
  are varied across cycles (generic terms, method-specific terms such as
  ``NEVPT2 zero-field splitting,'' and metal-specific terms) to broaden
  coverage.
\item \textbf{Fetch.} PDFs are downloaded \emph{only} from the open-access
  locations reported by the API (\status{fetched}, or \status{fetch\_failed}
  with a reason). Paywalled content and third-party mirrors are never
  accessed.
\item \textbf{Text extraction.} Each PDF is converted to plain text
  (\status{text\_ready}). If the text is garbled beyond reliable reading
  (equations or tables shredded), the paper is marked
  \status{text\_unreadable} and abandoned---numbers are never guessed from
  corrupted output.
\item \textbf{Read and judge.} The agent reads the extracted text and
  applies the scope test above. Papers that fail it are marked
  \status{not\_usable} with a logged reason.
\item \textbf{Extract rows.} For each distinct compound~$\times$~method
  result in a usable paper, the agent builds one schema-conforming record
  and appends it (\status{extracted}, or \status{extracted\_partial} when
  key fields were unavailable but the record is still useful).
\item \textbf{Log.} Every status change appends a line
  (\col{timestamp, key, doi, action, result, reasoning}) to an append-only
  audit log, written so that a chemist can check the reasoning after the
  fact.
\end{enumerate}

\paragraph{Supporting information.}
Computational details---active-space composition, geometries, software, and
state energies---are frequently relegated to the Supporting Information.
When the main text defers to the SI or leaves a key field unstated, the SI
is retrieved (figshare DOIs first, then PMC open-access packages) and
text-extracted before any field is declared ``not stated.'' Atomic
coordinates found in the SI are stored \emph{verbatim} with an explicit
provenance line and take precedence over any computationally generated
geometry.

\section{Record schema}

Each row of \col{literature.csv} is one compound~$\times$~method result. The
32 columns are grouped below by role. Multiple compounds from a single paper
receive suffixed identifiers (\col{-a}, \col{-b}, \dots); \col{reference_doi}
is mandatory and a row without it cannot be added.

\begingroup
\small
\begin{longtable}{@{}p{0.30\textwidth}p{0.63\textwidth}@{}}
\toprule
\textbf{Column} & \textbf{Meaning} \\
\midrule
\endfirsthead
\toprule
\textbf{Column} & \textbf{Meaning} \\
\midrule
\endhead
\midrule \multicolumn{2}{r}{\emph{continued on next page}} \\
\endfoot
\bottomrule
\endlastfoot

\multicolumn{2}{@{}l}{\textit{Identity and structure}}\\
\col{entry_id}        & Record key: the paper key plus \col{-a}/\col{-b}/\dots{} per compound. \\
\col{structure_file}  & Filename of a saved geometry, if any (else empty). \\
\col{compound_name}   & Chemical name / description of the species. \\
\col{formula}         & Molecular formula, as stated. \\
\col{metal_center}    & Metal element(s); \col{none} for metal-free organics. \\
\col{metal_ox_state}  & Formal oxidation state of the metal. \\
\col{d_electron_count}& $d^n$ count (arithmetic from the oxidation state). \\
\col{ligand_set}      & Coordinating ligands / donor set. \\
\col{point_group}     & Molecular point-group symmetry. \\
\col{geometry_source} & Origin of the geometry (optimized, X-ray, prior work, \dots). \\
\addlinespace
\multicolumn{2}{@{}l}{\textit{Method and level of theory}}\\
\col{method}          & Multireference method(s) used. \\
\col{software}        & Program package. \\
\col{basis_set}       & One-electron basis / auxiliary bases. \\
\col{relativistic_treatment} & Scalar-relativistic Hamiltonian (DKH, X2C, ECP, \dots). \\
\col{soc_included}    & Whether spin--orbit coupling was included. \\
\col{ss_included}     & Whether the spin--spin term was included. \\
\col{correlation_correction} & Post-CASSCF dynamic-correlation correction applied
  (NEVPT2, CASPT2, RASPT2, MRCI, \dots), or \col{none} for a variational
  CASSCF/RASSCF/DMRG-only result; derived from \col{method}. \\
\addlinespace
\multicolumn{2}{@{}l}{\textit{Active space}}\\
\col{active_space_nel}  & Number of active electrons. \\
\col{active_space_norb} & Number of active orbitals. \\
\col{active_space_orbital_description} & Chemical identity of the active orbitals. \\
\col{active_space_protocol} & The paper's active-space build-up sequence, quoted as
  ``(nel,norb) step; (nel,norb) step; \dots''. \\
\addlinespace
\multicolumn{2}{@{}l}{\textit{States and energetics}}\\
\col{multiplicities_studied} & Spin multiplicities considered. \\
\col{nroots_per_mult}   & Number of roots per multiplicity (state averaging). \\
\col{ground_state_mult} & Ground-state multiplicity. \\
\col{electronic_structure_description} & Ground-state term symbol (prefixed ``Ground term: \dots'' when present) and low-lying state energies (units stated; SOC-free vs.\ SOC-corrected noted). \\
\col{Other}             & Other reported quantities (gaps, ZFS, \dots). \\
\addlinespace
\multicolumn{2}{@{}l}{\textit{Provenance}}\\
\col{reference_short}   & Short citation. \\
\col{reference_doi}     & DOI (\textbf{mandatory}). \\
\col{year}              & Publication year. \\
\col{notes}             & Extraction caveats, unit conversions, and assumptions. \\
\col{entry_type}        & How the entry was produced: \col{llm_mining} (LLM-mined from the literature, the current default), \col{llm_reproduced} (a calculation an LLM re-ran), or \col{human} (human-entered). \\
\col{mining_model}      & Identifier of the LLM (Claude model) that performed the extraction, stamped at insertion time. \\
\end{longtable}
\endgroup

\section{Extraction principles}

The value of the database rests on a single discipline: \textbf{copy, never
infer.} The following rules are applied without exception.

\begin{description}
\item[Every value is stated in the paper.] A stored value must appear in the
  paper's text, tables, or SI, or be an exact arithmetic restatement of
  something that does (e.g.\ converting cm$^{-1}$ to eV). The conversion is
  then noted in \col{notes}. Values are never estimated or filled in from
  general chemical knowledge.
\item[Empty wins.] A field the paper does not state is left \emph{empty}. An
  empty cell is a deliberate, honest signal of absence and is always
  preferred over a plausible guess.
\item[Deliberate absence vs.\ unknown.] For metal-free organic compounds,
  \col{metal_center} is set to the literal \col{none}---a positive assertion
  that there is no metal---while \col{metal_ox_state}, \col{d_electron_count},
  and \col{ligand_set} are left empty. This distinguishes ``no metal
  exists'' from ``a metal exists but its property was not stated.''
\item[Sanctioned arithmetic.] \col{d_electron_count} may be computed from a
  stated oxidation state (it is arithmetic); if the paper does not state the
  oxidation state explicitly, the inference is flagged in \col{notes}. Unit
  conversions are likewise permitted and annotated.
\item[One row per compound, condensed.] A paper reporting many tables for
  one species yields a single, condensed record, not one row per table.
  Distinct compounds---or genuinely distinct method results---from the same
  paper get separate suffixed rows.
\item[Protocol quoted faithfully.] \col{active_space_protocol} is populated
  only when the paper describes its own active-space build-up, and is quoted
  in the paper's own terms; this column is the one the downstream decision
  agent reads most closely.
\item[Append-only.] Existing rows are never edited or deleted; corrections
  are added to \col{notes}. Together with the audit log this makes the
  corpus fully reconstructible and reviewable.
\end{description}

\section{Current database statistics}

Table~\ref{tab:stats} summarizes the present snapshot. Coverage percentages
reflect the ``empty wins'' discipline: fields such as \col{formula} and
\col{electronic_structure_description} are populated only where the source
paper stated them, so partial coverage is expected and informative rather
than a defect.

\begin{table}[h]
\centering
\small
\begin{tabular}{@{}lr@{}}
\toprule
\textbf{Quantity} & \textbf{Value} \\
\midrule
Calculation records (rows)            & 667 \\
Distinct publications (DOIs)          & 403 \\
Publication-year range                & 1985--2026 \\
Metal-free organic records (\col{metal_center}=\col{none}) & 408 \\
Distinct elements represented         & 58 \\
\midrule
\multicolumn{2}{@{}l}{\textit{Field coverage}}\\
\col{compound_name}                   & 667 (100\%) \\
\col{metal_center}                    & 667 (100\%) \\
\col{reference_doi}                   & 667 (100\%) \\
\col{year}                            & 667 (100\%) \\
\col{active_space_protocol}           & 640 (96\%) \\
\col{active_space_norb}               & 601 (90\%) \\
\col{active_space_nel}                & 589 (88\%) \\
\col{basis_set}                       & 437 (66\%) \\
\col{formula}                         & 293 (44\%) \\
\col{electronic_structure_description} & 259 (39\%) \\
\midrule
\multicolumn{2}{@{}l}{\textit{Method family (by leading method named)}}\\
CASSCF                                & 357 \\
CASPT2                                & 181 \\
NEVPT2                                & 98 \\
RASSCF                                & 10 \\
DMRG                                  & 8 \\
MRCI                                  & 7 \\
RASPT2                                & 1 \\
\midrule
\multicolumn{2}{@{}l}{\textit{Correlation correction applied (\col{correlation_correction})}}\\
\col{none} (variational CASSCF/RASSCF/DMRG only) & 338 \\
CASPT2                                & 208 \\
NEVPT2                                & 98 \\
MRCI                                  & 62 \\
CIPSI / DDCI / XMCQDPT2 / RASPT2 / other & 31 \\
\midrule
\multicolumn{2}{@{}l}{\textit{Entry provenance (\col{entry_type})}}\\
\col{llm_mining} (all records, \col{claude-opus-4-8}) & 667 \\
\bottomrule
\end{tabular}
\caption{Database snapshot statistics. Method-family counts classify each
record by the highest-level method named in \col{method}; many records also
report a lower-level reference calculation. Correlation-correction counts are
per token, so a record reporting several corrections (e.g.\ \col{CASPT2; MRCI})
contributes to each.}
\label{tab:stats}
\end{table}

\section{Provenance, auditability, and limitations}

\paragraph{Auditability.}
Three artifacts make the corpus reviewable. \col{papers_index.csv} records
the status of every candidate paper (including those rejected and why);
\col{logs/extractions.csv} is an append-only, timestamped narrative of every
decision written for a human chemist to check; and each record's \col{notes}
field carries the per-value caveats. Copyrighted artifacts---downloaded PDFs
and extracted text---are never committed to version control.

\paragraph{Limitations.}
The corpus is bounded by what is \emph{open-access} and \emph{machine-%
readable}: paywalled papers are absent by design, and papers whose PDFs
extract to garbled text are dropped rather than risk transcription errors.
Because the pipeline is deliberately conservative---empty cells over
guesses---absence of a value in the database means ``not reliably extracted,''
not ``not true.'' Extraction judgments, while logged, are made by a
language-model agent and carry a nonzero error rate; the audit log exists
precisely so that any record can be traced to its source and checked.

\paragraph{Browsing.}
The database ships with a self-contained static web interface
(\col{site/index.html}) that provides full-text search and faceted filtering
by compound name, chemical formula, and element, with each record expandable
to all 32 fields and linked to its DOI. It is regenerated from the CSV by
\col{site/build_site.py}. The site also has a \emph{Contribute} form through
which community members can submit an entry; submissions are written to a
separate review queue (\col{database/contributions.csv}) and are never merged
automatically---the maintainer vets each one against its DOI before it is
promoted into the database.

\end{document}
